\begin{summarybox}
\textbf{\textcolor{CSummary}{一句话核心}}
\textbf{椭圆上的动点与\emph{长轴两端点}连线，两斜率的乘积不随动点改变，且一定为负。}

\medskip
\textbf{\textcolor{CSummary}{双方向公式（$a>b>0$ 均表示长短半轴）}}
\[
 \boxed{\begin{aligned}
 &\frac{x^2}{a^2}+\frac{y^2}{b^2}=1
 &&\Longrightarrow\quad k_{PA_1}k_{PA_2}=-\frac{b^2}{a^2}=e^2-1,\\
 &\frac{y^2}{a^2}+\frac{x^2}{b^2}=1
 &&\Longrightarrow\quad k_{PB_1}k_{PB_2}=-\frac{a^2}{b^2}.
 \end{aligned}}
\]
这里 $A_{1,2}=(\mp a,0)$，$B_{1,2}=(0,\mp a)$；$P$ 必须在椭圆上且异于所选两端点。

\medskip
\textbf{\textcolor{CSummary}{3 条记忆路线}}
\begin{itemize}[leftmargin=*,itemsep=0.35em,topsep=0.3em]
 \item \textbf{最快：}斜率积 $=\frac{y^2}{x^2-a^2}$；椭圆方程给出 $y^2=-\frac{b^2}{a^2}(x^2-a^2)$。
 \item \textbf{最直观：}圆上直径所对的两直线垂直，斜率积 $-1$；把圆竖向压缩 $b/a$，两斜率各乘 $b/a$。
 \item \textbf{最稳妥：}先看长轴方向，再看端点与点的位置，最后计算斜率或用定值。
\end{itemize}

\medskip
\textbf{\textcolor{CSummary}{圆 $\longleftrightarrow$ 椭圆：5 秒回忆链}}
\[
 \underbrace{u^2+v^2=a^2}_{\text{圆}}
 \quad\xrightarrow[\text{横坐标固定}]{\text{纵坐标乘 }b/a}\quad
 \underbrace{\frac{x^2}{a^2}+\frac{y^2}{b^2}=1}_{\text{椭圆}}.
\]
\[
 \boxed{Q(u,v)\mapsto P(u,\tfrac ba v)},\qquad
 \boxed{P(x,y)\mapsto Q(x,\tfrac ab y)}.
\]
\textbf{为什么是 $b/a$：}两条曲线在同一横坐标处满足
$\frac{v^2}{a^2}=1-\frac{u^2}{a^2}=\frac{y^2}{b^2}$，
同侧取同号，便有 $y/b=v/a$。
\textbf{不变：}横坐标、垂足 $H$、整条横轴（含 $A_1,O,A_2$）、归一化高度 $v/a=y/b$；
\textbf{变化：}每条斜率乘 $b/a$，角度通常不保留。
\textbf{圆里先有 $k_1k_2=-1$，椭圆里才有 $k'_1k'_2=-(b/a)^2$。}

\medskip
\textbf{\textcolor{CSummary}{高考识别信号与自检}}
动点 + 椭圆 + 两个长轴顶点 + 斜率积 $\Rightarrow$ 立即想到第三定义。
自检两问：\textbf{第一，端点是不是长轴顶点？第二，点 $P$ 是否排除了端点？}

\medskip
\textbf{\textcolor{CSummary}{一道极简变式}}
若水平长轴的椭圆满足 $k_{PA_1}k_{PA_2}=-\frac9{25}$，
则 $e=\sqrt{1-\frac9{25}}=\boxed{\frac45}$。

\medskip
\textbf{\textcolor{CSummary}{逆向补充}}
对固定的水平两点 $(-a,0),(a,0)$，斜率积 $-\lambda$ 对应的曲线为
$\frac{x^2}{a^2}+\frac{y^2}{\lambda a^2}=1$（不包含这两个点）；
$0<\lambda<1$ 才说明指定两点是长轴端点。
\end{summarybox}
